% =====================================================================
% LAMPIRAN B: CONTOH SAMPEL MULTIMODAL TAMBAHAN
% =====================================================================
\chapter{Contoh Sampel Multimodal Tambahan}
\label{app:sample-multimodal}

Lampiran ini menyajikan dua contoh baris dataset tambahan yang melengkapi Gambar~\ref{fig:sample-row-1} pada Bagian~\ref{sec:dataset}. Kedua contoh berikut dipilih untuk menggambarkan variasi panjang narasi dan jenis kejadian, sehingga sifat multimodal serta komplementer antara citra dan teks dapat ditinjau lebih luas tanpa membebani pembahasan utama Bab~\ref{ch:metode}.

\begin{figure}[htbp]
    \centering
    \includegraphics[width=0.95\textwidth]{sample_row_2.png}
    \caption{Contoh 2 struktur data multimodal: foto saluran air mampet dengan narasi panjang yang menjelaskan dampak banjir terhadap warga sekitar, dipetakan ke label \textit{Kelurahan}. Narasi yang panjang menyediakan konteks tekstual kuat, sedangkan citra memvalidasi kondisi fisik saluran.}
    \label{fig:sample-row-2}
\end{figure}

\begin{figure}[htbp]
    \centering
    \includegraphics[width=0.95\textwidth]{sample_row_3.png}
    \caption{Contoh 3 struktur data multimodal: foto pembatas jalan dengan cat yang sudah pudar, narasi permohonan pengecatan ulang oleh warga, dipetakan ke label \textit{Dinas Perhubungan}. Kondisi visual pudar pada citra menjadi sinyal utama, sedangkan narasi menegaskan jenis tindakan yang diharapkan.}
    \label{fig:sample-row-3}
\end{figure}
